\section{Jurisdiction, Document Structure, and Multilingual Legal AI}
\label{sec:jurisdiction}

\subsection{Why legal information retrieval is jurisdiction-dependent}

Legal information retrieval differs from general-purpose information retrieval because the relevance of a
document depends not only on its subject matter but also on its legal authority, jurisdiction, temporal
validity, and relationship to the question being asked. Two passages may discuss the same legal
principle while having different implications because they originate from different jurisdictions or
occupy different positions within a legal system.

For example, a question concerning criminal procedure in Pakistan should prioritize applicable Pakistani
legislation and relevant Pakistani judicial decisions. A semantically similar provision from another
jurisdiction may help explain a general concept, but it should not automatically be treated as governing
authority. Jurisdiction-sensitive retrieval must therefore consider several factors:
\begin{itemize}[leftmargin=1.4em, itemsep=2pt, topsep=2pt]
\item \textbf{Territorial and institutional jurisdiction:} whether the authority applies to the relevant
court, territory, or legal institution.
\item \textbf{Authority type:} whether a document is legislation, a regulation, a binding judgment, a
persuasive judgment, or secondary commentary.
\item \textbf{Temporal validity:} whether the provision or interpretation was applicable at the relevant
time.
\item \textbf{Procedural context:} whether the authority addresses the procedural stage and legal issue
involved.
\item \textbf{Precedential relationships:} whether a judgment follows, distinguishes, overrules, or
otherwise relates to another decision.
\end{itemize}

These considerations expose a limitation of conventional semantic retrieval. An embedding model may
identify two passages as highly similar without determining whether either passage is legally
applicable. Consequently, jurisdictional and temporal constraints should influence retrieval and
verification, not merely appear as explanatory text in the final answer.

\subsection{Legal hierarchy and document structure}

Legal documents contain multiple layers of organization. Statutes may be divided into parts, chapters,
sections, subsections, clauses, and provisos. Judgments may contain procedural histories, questions of
law, arguments, judicial reasoning, and final orders. These structural elements help determine the
meaning and applicability of individual passages.

A conventional chunking strategy often divides documents into fixed-size text segments or overlapping
windows. Although computationally convenient, this can separate a provision from its heading, definition,
exception, proviso, or parent section. The resulting fragment may remain semantically understandable
while losing context that is essential to legal interpretation. Consider a hypothetical example:

\begin{quote}
\textit{Illustrative statutory structure (schematic):}
\begin{itemize}[leftmargin=1.3em, itemsep=1pt, topsep=2pt]
\item \textbf{Chapter III: General Conditions}
  \begin{itemize}[leftmargin=1.2em, itemsep=1pt, topsep=1pt]
  \item \textbf{Section 12 (general rule):} a person may exercise the stated right subject to the
  conditions specified in this section.
    \begin{itemize}[leftmargin=1.2em, itemsep=1pt, topsep=1pt]
    \item \textbf{Subsection (2) (limitation):} the right is subject to the additional restrictions
    specified in the subsection.
    \item \textbf{Proviso:} an exception applies where the specified conditions are satisfied.
    \end{itemize}
  \end{itemize}
\end{itemize}
This is an illustrative structure only, not an actual statutory provision.
\end{quote}

If retrieval returns only the general rule, the generated answer may omit the limitation or exception.
If it returns only the proviso, the system may fail to explain the general rule to which that exception
applies. Structure-aware retrieval addresses this problem by preserving the relationship between a
passage and its parent document elements. Useful techniques include:
\begin{enumerate}[leftmargin=1.4em, itemsep=2pt, topsep=2pt]
\item \textbf{Hierarchical metadata:} storing the document, chapter, section, subsection, and paragraph
identifiers associated with each passage.
\item \textbf{Context-preserving chunking:} retaining headings, definitions, exceptions, and relevant
neighboring passages.
\item \textbf{Parent-child retrieval:} retrieving a focused passage while allowing the system to recover
its broader context.
\item \textbf{Reference resolution:} linking references to other sections, definitions, schedules, and
related legal instruments.
\item \textbf{Provenance preservation:} maintaining the source document, version, location, and other
information needed to verify the passage.
\end{enumerate}

Structure-aware indexing and knowledge-graph retrieval are complementary. Hierarchical indexing preserves
the organization of a document, whereas a knowledge graph represents explicit relationships between
entities, legal provisions, judgments, and other objects. A system may use both, but each requires
separate construction and evaluation.

\subsection{Pakistani legal information retrieval}

Pakistan provides a useful case study of the challenges involved in building jurisdiction-aware legal
AI. Its legal information needs span constitutional provisions, statutes, procedural codes, judicial
decisions, and other legal materials. A useful system must preserve the relationships among these
sources rather than treating every retrieved passage as an interchangeable text fragment.

For example, the Constitution of Pakistan, the Pakistan Penal Code (PPC), the Code of Criminal Procedure
(CrPC), the Qanun-e-Shahadat Order (QSO), and judicial decisions serve different legal functions. The
appropriate source depends on the question. A constitutional question may require a constitutional
provision and relevant case law, while a question about a criminal procedure may require the applicable
procedural provision and related judicial interpretation.

Three challenges are particularly important. First, \textit{document coverage and authority.} A corpus
may contain a large volume of legal text but still omit relevant amendments, judgments, or authoritative
documents. Corpus size alone is therefore an inadequate measure of legal coverage. Second, \textit{legal
relationships.} A statutory provision may depend on a definition elsewhere in the same enactment or be
interpreted by a later judgment. Retrieval should preserve these relationships where relevant, rather
than relying exclusively on textual similarity. Third, \textit{temporal and jurisdictional
applicability.} A provision's wording, legal status, or interpretation may change over time. Systems
should record source dates and applicable versions, and distinguish the date of a judgment from the
period in which the underlying law applied.

A jurisdiction-aware Pakistani legal retrieval system should consequently combine source provenance,
hierarchical indexing, authority-aware ranking, and temporal metadata. These features can improve the
reliability of evidence selection, although they do not eliminate the need for legal interpretation and
claim-level verification.

\subsection{English, Urdu, and low-resource legal retrieval}

Multilingual legal AI introduces an additional challenge: the language of a user's query may differ from
the language, terminology, or script of the authoritative source. This is particularly relevant when
users express legal questions in Urdu or Roman Urdu while the source material they need is written in
English.

A multilingual system must bridge differences in vocabulary, spelling, transliteration, morphology, and
legal terminology. Direct translation alone may be insufficient because legal concepts can have
jurisdiction-specific meanings that do not map exactly across languages. Four technical approaches are
particularly relevant, as summarized in Table~\ref{tab:multilingual}.

\begin{table*}[tp]
\centering
\caption{Technical approaches to multilingual legal retrieval.}
\label{tab:multilingual}
\footnotesize
\setlength{\tabcolsep}{5pt}
\renewcommand{\arraystretch}{1.3}
\begin{tabularx}{\textwidth}{>{\raggedright\arraybackslash}p{3.2cm} >{\raggedright\arraybackslash}X >{\raggedright\arraybackslash}X >{\raggedright\arraybackslash}X}
\toprule
\textbf{Approach} & \textbf{How it works} & \textbf{Main advantage} & \textbf{Main limitation} \\
\midrule
Query translation & Translates the query into the language of the source corpus. & Enables use of existing monolingual retrieval infrastructure. & Translation may distort legal terms or procedural meaning. \\
Multilingual dense retrieval & Maps queries and documents from different languages into a shared embedding space. & Supports cross-language semantic matching. & Similar embeddings do not guarantee legal-concept equivalence. \\
Transliteration normalization & Normalizes alternative representations of words across scripts. & Helps retrieve Roman Urdu and Urdu variants. & Many transliterations are ambiguous or inconsistent. \\
Legal terminology alignment & Links expressions to jurisdiction-specific legal concepts and authoritative terms. & Improves conceptual consistency. & Requires curated terminology, entity linking, or legal-domain resources. \\
\bottomrule
\end{tabularx}
\end{table*}

A practical architecture can combine these methods rather than relying on one. Language identification
and transliteration normalization can first process the query. Legal entity linking and terminology
alignment can then identify important statutes, provisions, courts, or legal concepts. Retrieval can
search both the original query and its normalized or translated variants, followed by rank fusion and
legal-context filtering.

The system should preserve the original user query and the original authoritative passage. Translation
can support retrieval, but the final answer should not silently substitute a translated interpretation
for the governing source text.

Evaluation should also be multilingual. An English-only benchmark cannot establish whether the system
works reliably for Urdu or Roman Urdu. Separate evaluation should measure retrieval recall, citation
support, legal correctness, and robustness to spelling and transliteration variation across each
supported query language.

\subsection{Integrating jurisdiction, hierarchy, and multilingual processing}

The three challenges discussed above are interdependent. Language normalization may help locate a
relevant provision, but document structure is needed to preserve its context. Hierarchical indexing may
identify the correct section, but jurisdictional and temporal checks are still required to determine its
applicability. A conceptual architecture connecting these stages is shown in Figure~\ref{fig:mlarch}.

% Conceptual architecture: multilingual, jurisdiction- and structure-aware legal retrieval. Vector TikZ.
\begin{figure}[t]
\centering
\resizebox{\columnwidth}{!}{%
\begin{tikzpicture}[
  font=\footnotesize,
  stage/.style={draw, thick, rounded corners=2pt, align=center, text width=6.6cm,
                minimum height=8mm, inner sep=3pt, fill=legalblue!6, draw=legalblue!70},
  sub/.style={draw, thick, rounded corners=2pt, align=center, text width=3.1cm,
              minimum height=8mm, inner sep=3pt, fill=black!4, draw=black!55},
  fl/.style={-{Latex[length=2.4mm]}, draw=black!50, line width=0.8pt}]

  \node[stage] (q) at (0,0) {\textbf{Multilingual legal query}\\\scriptsize English, Urdu, or Roman Urdu};
  \node[stage] (norm) at (0,-1.5) {\textbf{Query normalization and legal concept alignment}\\\scriptsize language identification, transliteration, terminology, entity linking};
  \node[sub] (idx) at (-1.85,-3.15) {\textbf{Hierarchical legal index}\\\scriptsize documents, sections, subsections, provenance};
  \node[sub] (kg) at ( 1.85,-3.15) {\textbf{Legal knowledge graph}\\\scriptsize entities, relationships, cross-references, authority links};
  \node[stage] (ret) at (0,-4.85) {\textbf{Hybrid retrieval and evidence ranking}\\\scriptsize lexical and dense retrieval, rank fusion, reranking, jurisdiction filters};
  \node[stage] (resp) at (0,-6.35) {\textbf{Verified legal response}\\\scriptsize source-grounded claims, citation checks, temporal validation, uncertainty handling};

  \draw[fl] (q) -- (norm);
  \draw[fl] (norm.south) -- (idx.north);
  \draw[fl] (norm.south) -- (kg.north);
  \draw[fl] (idx.south) -- (ret.north);
  \draw[fl] (kg.south) -- (ret.north);
  \draw[fl] (ret) -- (resp);
\end{tikzpicture}}
\caption{Conceptual architecture connecting multilingual query understanding, legal document structure,
knowledge graphs, retrieval, and evidence verification. Synthesized for this survey; its individual
components require empirical evaluation before any performance improvement can be claimed.}
\label{fig:mlarch}
\end{figure}


This architecture offers a research framework for connecting multilingual query understanding, legal
document structure, knowledge graphs, retrieval, and evidence verification. It also provides a basis for
modular implementation: individual components can be replaced or evaluated independently without
redesigning the entire platform.

\subsection{Section synthesis}

Jurisdiction, document structure, and language are not peripheral concerns in legal information
retrieval. They determine which evidence should be retrieved, how its meaning should be interpreted, and
whether it applies to the question being asked.

For Pakistani legal AI, the main implication is that effective retrieval requires more than semantic
similarity. Systems should preserve legal hierarchy, authority relationships, temporal metadata, source
provenance, and cross-language concept alignment. These requirements lead directly to the open
challenges discussed next.
